%%% Document Formatting
\documentclass[11pt,fleqn]{article}
\usepackage[a4paper,
            bindingoffset=0.2in,
            left=0.75in,
            right=0.75in,
            top=0.8in,
            bottom=0.8in,
            footskip=.25in]{geometry}
\setlength\parindent{10pt}

%%% Imports
\usepackage{amsmath, amsfonts, amssymb, amsthm}
\usepackage{dsfont}
\numberwithin{equation}{section}
\usepackage{mathtools}
\usepackage{cancel}
\usepackage{empheq}
\usepackage{algorithm}
\usepackage{algpseudocode}  % preamble

% Visualization
\usepackage{graphicx}


%%% Formatting
\usepackage{hyperref}
\hypersetup{
    colorlinks=true,
    linkcolor=blue,
    filecolor=magenta,      
    urlcolor=cyan,
    pdftitle={SCSI Lecture Notes},
    pdfpagemode=FullScreen,
}
\urlstyle{same}

%%% Custom Environments
 

% theorems
% Dependencies
\usepackage[usenames,dvipsnames]{xcolor}
\usepackage{amsmath, amsfonts, mathtools, amsthm, amssymb}
\usepackage{thmtools}
\usepackage[framemethod=TikZ]{mdframed}
\mdfsetup{skipabove=1em,skipbelow=0em}

\theoremstyle{definition}


\makeatother
\usepackage{thmtools}
\usepackage[framemethod=TikZ]{mdframed}
\mdfsetup{skipabove=1em,skipbelow=0em}


\theoremstyle{definition}

\declaretheoremstyle[
    headfont=\bfseries\color{ForestGreen!70!black}, bodyfont=\normalfont,
    mdframed={
        linewidth=2pt,
        rightline=false, topline=false, bottomline=false,
        linecolor=ForestGreen, backgroundcolor=ForestGreen!5,
    }
]{thmgreenbox}

\declaretheoremstyle[
    headfont=\bfseries\color{NavyBlue!70!black}, bodyfont=\normalfont,
    mdframed={
        linewidth=2pt,
        rightline=false, topline=false, bottomline=false,
        linecolor=NavyBlue, backgroundcolor=NavyBlue!5,
    }
]{thmbluebox}

\declaretheoremstyle[
    headfont=\bfseries\color{NavyBlue!70!black}, bodyfont=\normalfont,
    mdframed={
        linewidth=2pt,
        rightline=false, topline=false, bottomline=false,
        linecolor=NavyBlue
    }
]{thmblueline}

\declaretheoremstyle[
    headfont=\bfseries\color{RawSienna!70!black}, bodyfont=\normalfont,
    mdframed={
        linewidth=2pt,
        rightline=false, topline=false, bottomline=false,
        linecolor=RawSienna, backgroundcolor=RawSienna!5,
    }
]{thmredbox}

\declaretheoremstyle[
    headfont=\bfseries\color{RawSienna!70!black}, bodyfont=\normalfont,
    numbered=no,
    mdframed={
        linewidth=2pt,
        rightline=false, topline=false, bottomline=false,
        linecolor=RawSienna, backgroundcolor=RawSienna!1,
    },
    qed=\qedsymbol
]{thmproofbox}

\declaretheoremstyle[
    headfont=\bfseries\color{NavyBlue!70!black}, bodyfont=\normalfont,
    numbered=no,
    mdframed={
        linewidth=2pt,
        rightline=false, topline=false, bottomline=false,
        linecolor=NavyBlue, backgroundcolor=NavyBlue!1,
    },
]{thmexplanationbox}

\declaretheoremstyle[
    headfont=\bfseries\color{Plum!70!black}, bodyfont=\normalfont,
    mdframed={
        linewidth=2pt,
        rightline=false, topline=false, bottomline=false,
        linecolor=Plum, backgroundcolor=Plum!5,
    }
]{thmpurplebox}

\declaretheoremstyle[
    headfont=\bfseries\color{Sepia!70!black}, bodyfont=\normalfont,
    mdframed={
        linewidth=2pt,
        rightline=false, topline=false, bottomline=false,
        linecolor=Sepia, backgroundcolor=Sepia!5,
    }
]{thmsepiabox}

\declaretheoremstyle[
    headfont=\bfseries\color{ForestGreen!70!black}, bodyfont=\normalfont,
    mdframed={
        linewidth=2pt,
        rightline=false, topline=false, bottomline=false,
        linecolor=ForestGreen
    }
]{thmgreenline}

% \declaretheoremstyle[headfont=\bfseries, bodyfont=\normalfont, mdframed={ nobreak } ]{thmgreenbox}
% \declaretheoremstyle[headfont=\bfseries, bodyfont=\normalfont, mdframed={ nobreak } ]{thmredbox}
% \declaretheoremstyle[headfont=\bfseries, bodyfont=\normalfont]{thmbluebox}
% \declaretheoremstyle[headfont=\bfseries, bodyfont=\normalfont]{thmblueline}
% \declaretheoremstyle[headfont=\bfseries, bodyfont=\normalfont, numbered=no, mdframed={ rightline=false, topline=false, bottomline=false, }, qed=\qedsymbol ]{thmproofbox}
% \declaretheoremstyle[headfont=\bfseries, bodyfont=\normalfont, numbered=no, mdframed={ nobreak, rightline=false, topline=false, bottomline=false } ]{thmexplanationbox}

\declaretheorem[style=thmgreenbox, name=Definition]{definition}
\declaretheorem[style=thmgreenbox, name=Assumption]{assumption}
\declaretheorem[style=thmbluebox, numbered=no, name=Example]{eg}
\declaretheorem[style=thmredbox, name=Proposition]{prop}
\declaretheorem[style=thmredbox, name=Theorem]{theorem}
\declaretheorem[style=thmredbox, name=Lemma]{lemma}
\declaretheorem[style=thmredbox, numbered=no, name=Corollary]{corollary}
\declaretheorem[style=thmpurplebox, name=Problem]{problem}
\declaretheorem[style=thmsepiabox, numbered=no, name=Fact]{fact}
\declaretheorem[style=thmgreenline, name=Property]{property}

\declaretheorem[style=thmproofbox, name=Proof]{replacementproof}
\renewenvironment{proof}[1][\proofname]{\vspace{-10pt}\begin{replacementproof}}{\end{replacementproof}}


\declaretheorem[style=thmexplanationbox, name=Proof]{tmpexplanation}
\newenvironment{explanation}[1][]{\vspace{-10pt}\begin{tmpexplanation}}{\end{tmpexplanation}}

\declaretheorem[style=thmblueline, numbered=no, name=Remark]{remark}
\declaretheorem[style=thmblueline, numbered=no, name=Note]{note}

\newmdenv[ 
  topline=false,
  bottomline=false,
  skipabove=\topsep,
  skipbelow=\topsep
]{sidework}

\title{Homework \#2}
\author{Clark Miyamoto (cm6627@nyu.edu)}
\date{}

\begin{document}

\maketitle

\subsection*{Question 1)} Let $u,u' \overset{iid}{\sim} q$, and consider probability measures $\pi = \delta_u$ and $\pi' = \delta_{u'}$. Show 
\begin{align}
	d_{rand}(\pi, \pi')^2 = 2 \sup_{|f|_{\infty} \leq 1} \text{Var}_q[f]
\end{align}

\paragraph{Solution)} Recall
\begin{align}
	d_{rand}(\pi, \pi') = \sup_{|f|_{\infty} \leq 1} \left |  \mathbb E \left[\left( \mathbb E_\pi[f] - \mathbb E_{\pi'}[f] \right)^2 \right]\right  |^{1/2}
\end{align}
Since $\pi = \delta_u$, the expectation evaluates to $\mathbb E_\pi[f] = \int f(x) \, \delta_u(x) \, dx = f(u)$
\begin{align}
	d_{rand}(\pi, \pi')= \sup_{|f|_{\infty} \leq 1} \left | \mathbb E_{u,u' \sim q} \left[\left( f(u) - f(u' \right)^2 \right]\right|^{1/2}
\end{align}
Now notice the expectation inside has a nice  expression
\begin{align}
& \mathbb E_{u,u' \sim p}[(f(u) - f(u'))^2]\\
& = \int (f(u) - f(u'))^2 \, p(u) \, p'(u') \, du \, du'\\
& = \int [f(u)^2 + f(u')^2 - 2 f(u) f(u')] p(u) p'(u') \, du \, du'\\
& = \mathbb E_p[f^2] + \mathbb E_p[f^2] - 2 \mathbb E_p[f]^2\\
& = 2 \, \text{Var}_p[f]
\end{align}
So the metric over random probability measures becomes
\begin{align}
	\mathcal D := d_{rand}(\pi, \pi') = 2 \sup_{|f|_{\infty} \leq 1}  \sqrt{ \text{Var}_p[f] }
\end{align}
Notice $\phi(x) = \sqrt{x}$ is a monotonic function on $x \in [0, \infty)$. I think $\sup$ commutes with monotonic functions (assuming $\sup$ exists). Therefore $\mathcal D = 2 \left( \sup_{|f|_{\infty }} \text{Var}_p[f] \right)^{1/2}$, squaring both sides concludes the proof.


\subsection*{Question 3)} 
The optimization problem
\begin{align}
	\psi^* & \in \arg \min_{\psi \in \mathcal K} J^N(\psi) \label{eqn:opt}\\
	J^N(\psi) & = \frac{1}{2} \sum_{n=1}^N |y^{(n)} - \psi(u^{(n)}) |^2 + \frac{\lambda}{2} \|\psi\|_{\mathcal K}^2
\end{align}
where $\mathcal K$ is an RKHS with kernel $c(\cdot, \cdot)$. We defined $\langle f, g \rangle_{\mathcal K} = \sum_{i=1}^\infty \frac{\langle f, \varphi_i\rangle_{L^2} \langle g , \varphi_i \rangle_{L^2}}{\sigma_i}$.  Recall that the resulting function has the form $\psi^*(u) = \sum_{n=1}^N \alpha_n^* c(u, u^{(n)})$. Show the coefficients $\alpha^*$ solve
\begin{align}
	\alpha^* & \in \arg \min_{\alpha \in \mathbb R^N} A^N(\alpha)\\
	A^N(\alpha) & = \frac{1}{2N} \sum_{r=1}^N |y^{(r)} - \sum_{n=1}^N \alpha_n c(u^{(r)}, u^{(n)})|^2 + \frac{\lambda}{2N} \sum_{n,r=1}^N \alpha_n \alpha_r c(u^{(r)} , u^{(n)})
\end{align}
Deduce that finding $\alpha^*$ requires an inversion of a Gram matrix evaluate on the data.
\paragraph{Background}
\begin{theorem}
	[Mercer's Theorem] For integral operator $(C \varphi)(u) = \int c(u,u') \varphi(u') \, du'$, with kernel $c$; the kernel can be represent using the orthonormal basis $\left\{ \varphi_i \right\}_{i=1}^\infty$ (w.r.t. $L^2$ norm)
	\begin{align}
		c(u,u') = \sum_{i=1}^N \sigma_i \varphi_i(u) \, \varphi_i(u')
	\end{align}
\end{theorem}


\paragraph{Solution Part 1)} 
Since the solution to \ref{eqn:opt} has the form $\psi^*(u) = \sum_{n=1}^N \alpha_n^* c(u, u^{(n)})$. So we can plug this into the functional form of the loss function $J^N$.\\
\\
The only thing we need to compute is
\begin{align}
	\| \psi^*\|_{\mathcal K}^2 & = \langle \psi^*, \psi^* \rangle_{\mathcal K} \\
	& = \sum_{i=1}^\infty \frac{\langle \psi^*, \varphi_i \rangle_{L^2} \langle \psi^*, \varphi_i \rangle_{L^2}}{\sigma_i}\\
	&  = \sum_{i=1}^\infty \sum_{n,n'=1}^\infty  \frac{1}{\sigma_i} \alpha_n^* \alpha_{n'}^* \langle c(\cdot, u^{(n)}), \varphi_i \rangle_{L^2} \, \langle c(\cdot, u^{(n')} , \varphi_i\rangle_{L^2}
\end{align}
From Mercer's theorem, $c(u, u') = \sum_{i=1}^\infty \sigma_i \varphi_i (u) \varphi_i(u')$, so the inner product $\langle c(\cdot, x), \varphi_k\rangle_{L^2} = \sigma_k \varphi_k(x)$.
\begin{align}
	\| \psi^*\|_{\mathcal K}^2 & = \sum_{i=1}^\infty \sum_{n,n'=1}^N  \frac{1}{\sigma_i} \alpha_n^* \alpha_{n'}^* \sigma_i \varphi_i(u^{(n)}) \, \sigma_i \varphi_i(u^{(n')})\\
	& = \sum_{n,n'=1}^N \alpha_n^* \alpha_{n'}^*\sum_{i=1}^\infty \sigma_i \varphi_i(u^{(n)}) \varphi_i(u^{(n')})\\
	& = \sum_{n=1}^N\sum_{n'=1}^N \alpha_n^* \alpha_{n'}^* \, c(u^{(n)}, u^{(n')})
\end{align}

\paragraph{Solution Part 2)} 








\newpage
\appendix
\section*{Hilbert Spaces}
A Hilbert Space $\mathcal H$ is an vector space equiped with an inner product $\langle \cdot, \cdot \rangle : \mathcal H \times \mathcal H \to \mathbb R$. In particular the inner product should obey the following; let $x,y \in \mathcal H$
\begin{enumerate}
	\item $\langle y, x\rangle = \overline{\langle x, y \rangle}$
	\item Linear: $\langle ax_1 + bx_2, y\rangle = a \langle x_1 ,y \rangle + b \langle x_2, y\rangle$
	\item Positive definite: $\langle x, x \rangle \geq 0$ (if $x = 0$, then $\langle x, x \rangle = 0$.)
\end{enumerate} 
This inner product implicitly defines a norm $\|x \|_{\mathcal H} = \sqrt{\langle x, x \rangle}$.
































\end{document}